% Компиляция: xelatex (дважды для оглавления/ссылок).
\documentclass[11pt]{article}
\usepackage{fontspec}
\setmainfont{STIX Two Text}
\setmonofont{STIX Two Text}  % нет системного моно с кириллицей; texttt будет упрощённым
\usepackage{polyglossia}
\setdefaultlanguage{russian}
\setotherlanguage{english}
\usepackage{amsmath,amssymb,amsthm}
\usepackage[margin=1in]{geometry}
\usepackage{booktabs}
\usepackage{longtable}
\usepackage{microtype}
\usepackage[colorlinks=true,linkcolor=blue,urlcolor=blue]{hyperref}

\theoremstyle{plain}
\newtheorem{theorem}{Теорема}
\newtheorem{proposition}{Утверждение}
\newtheorem{corollary}{Следствие}
\theoremstyle{definition}
\newtheorem{assumption}{Предположение}
\newtheorem{definition}{Определение}

\newcommand{\lnx}{\ln x_2}
\newcommand{\dHfus}{\Delta H_{\mathrm{fus}}}
\newcommand{\Tm}{T_m}
\newcommand{\lng}{\ln\gamma_2}

\title{\textbf{Физико-информированное предсказание растворимости:\\
условие оптимальности физики и архитектура двойного заземления}\\[4pt]
\large Предварительный отчёт}
\author{Проект TGNN-Solv}
\date{\today}

\begin{document}
\maketitle

\begin{abstract}
\noindent
Мы исследуем, \emph{когда} явная термодинамическая декомпозиция помогает
предсказывать равновесную растворимость твёрдое--жидкость
$\lnx=-\Phi(T)-\lng$, а не просто \emph{помогает ли} она. Отчёт фиксирует
предварительные результаты по четырём направлениям: (A)~оценка неустранимого
шумового потолка меток; (B)~теория идентифицируемости и условие оптимальности
физической факторизации; (C)~архитектура \emph{двойного заземления}, где оба
фактора разложения (кристаллический член и активность через $\sigma$-профиль и
дифференцируемый COSMO-SAC) становятся одно-компонентными, внешне-супервизируемыми
объектами; (D)~оценка модели по утилите (ранжирование растворителей и калибровка
неопределённости), а не только по MAE. Ключевые количественные выводы:
неустранимый шумовой пол при честной межлабораторной оценке составляет
$\approx 0{,}3$--$0{,}7\ \log_{10}$, поэтому запас по абсолютной точности на новых
каркасах скромный (фактор ${\sim}2$, не порядок); на сопоставимом режиме
экстраполяции наша модель ($\approx 1{,}0\ \log_{10}$ RMSE) немного уступает
текущему SOTA FastSolv ($0{,}83$--$0{,}95$) и заметно опережает прежний SOTA
SolProp/Vermeire ($1{,}43$--$2{,}16$), причём SOTA --- это \emph{прямой чёрный
ящик}, а не физика; расщепление кристалл/активность строго неидентифицируемо из
растворимости и замыкается только внешними одно-компонентными данными; модель
ранжирует растворители заметно лучше ($\rho_{\text{Spearman}}=0{,}51$, top-1
$=0{,}48$), чем подсказывает её MAE, но её неопределённость сильно переуверена
(PICP$_{90}=0{,}13$).
\end{abstract}

\tableofcontents

\section{Введение и постановка задачи}
Цель проекта --- быстрый скрининг растворимости: для строки
$(\text{вещество }s,\ \text{растворитель }v,\ \text{температура }T)$ предсказать
$y=\lnx$ --- натуральный логарифм мольной доли растворённого вещества в
насыщенном растворе (типично $\lnx\in[-25,0]$). Центральный научный вопрос не
``может ли нейросеть предсказывать растворимость'', а \textbf{помогает ли явное
термодинамическое <<узкое место>>} по сравнению с тем же графовым энкодером,
обученным напрямую на $\lnx$ (контроль \textsc{DirectGNN}).

Накопленная диагностика проекта указывает, что прямое <<узкое место>> не даёт
выигрыша в точности (физический <<налог>> $\approx +0{,}09$ MAE на scaffold), а
основное ограничение точности --- \emph{представление} (RF на дескрипторах с
MAE $1{,}20$ опережает GNN). Настоящий отчёт переформулирует задачу: мы
доказываем \emph{условие}, при котором физика оптимальна, строим архитектуру,
которая это условие реализует конструктивно, и оцениваем модель метриками,
отражающими реальную задачу скрининга. Эта архитектура --- реализованная и
исправленная форма проектной линии \emph{ThermoCurve} (предсказание кривой,
аналитика вместо численного решателя, раздельная супервизия ветвей), которую мы
не отбросили, а довели до <<двойного заземления>> (\S\ref{sec:thermocurve}).

\section{Данные и неустранимый шумовой потолок}
\label{sec:noise}
Прежде чем улучшать модель, мы оценили, есть ли вообще запас по точности.
\textbf{Метод.} На сыром корпусе BigSolDB~v2.1 (атрибуция по источнику/DOI
сохранена; в обработанных сплитах она схлопнута) мы сгруппировали измерения по
$(\text{вещество},\text{растворитель},T{\pm}1\,\text{K})$ и среди групп с
$\geq 2$ независимыми источниками измерили внутригрупповое расхождение $\lnx$ ---
нижнюю границу неустранимого шума, который не может побить никакая модель.

\begin{table}[h]
\centering
\caption{Межисточниковое расхождение $\lnx$ (1687 источников, 4358
мульти-источниковых групп; лишь 97 --- точные дубликаты).}
\begin{tabular}{lrrrl}
\toprule
Вариант & медиана & среднее & p90 & CI$_{95}$(медианы) \\
\midrule
сырой & $0{,}031$ & $0{,}161$ & $0{,}450$ & $[0{,}029;\,0{,}033]$ \\
дедуп & $0{,}033$ & $0{,}165$ & $0{,}461$ & $[0{,}031;\,0{,}035]$ \\
\bottomrule
\end{tabular}
\end{table}

\noindent
По числу источников в группе расхождение растёт ($2\!\to\!0{,}027$;
$3\!\to\!0{,}091$; $4$--$5\!\to\!0{,}175$; $6{+}\!\to\!0{,}391$), но даже самые
перемеренные пары держатся много ниже ошибки модели.

\paragraph{Две оценки пола и честная сверка.} Числа выше --- в нат-логе мольной
доли; в $\log_{10}$ это $\approx 0{,}013$ (медиана) и $\approx 0{,}07$ (среднее).
Это \emph{оптимистичная нижняя} оценка: она измеряет согласие внутри одного
корпуса по совпадающей паре/$T$ и занижается ре-репортами и отбором совпадающих
точек. Более консервативную \emph{меж}лабораторную оценку даёт литература
(FastSolv, \S\ref{sec:external}): межлабораторное стандартное отклонение
$0{,}34\ \log_{10}$, типичные границы $0{,}5$--$0{,}7\ \log_{10}$. Истинный
неустранимый пол, по-видимому, лежит между этими оценками, ближе к
$0{,}3$--$0{,}7\ \log_{10}$ ($\approx 0{,}7$--$1{,}6$ в нат-логе).

\begin{quote}\textbf{Вывод (исправленный).} Точность \emph{выигрываема, но запас
скромный}: наш scaffold-RMSE $\approx 1{,}0\ \log_{10}$ выше реалистичного пола
$0{,}3$--$0{,}7\ \log_{10}$ примерно в $1{,}5$--$2$ раза, а текущий SOTA (FastSolv,
$0{,}83$--$0{,}95\ \log_{10}$ на экстраполяции) уже \emph{у предела алеаторной
неопределённости}. Значит резерв по абсолютной точности на новых каркасах ---
порядка фактора $2$, а не порядка величины; основной выигрыш надо искать в
ранжировании, температурной экстраполяции и интерпретируемости, а не в погоне за
MAE $\to 0$.\end{quote}

\section{Термодинамическая основа}
Для компонента $i$ в жидкой фазе $\mu_i^{\text{ж}}=\mu_i^{0}+RT\ln(\gamma_i x_i)$.
Из равенства химпотенциалов твёрдой и жидкой фаз растворённого вещества (индекс
$2$) следует основное уравнение равновесия твёрдое--жидкость (SLE):
\begin{equation}
\lnx = -\underbrace{\frac{\dHfus}{R}\!\Big(\frac1T-\frac1\Tm\Big)}
            _{\Phi(T)\ \text{(кристаллический член, зависит от }s)}
       \;-\;\underbrace{\lng(x_2,T)}_{\text{активность, зависит от }(s,v)},
\label{eq:sle}
\end{equation}
(в приближении $\Delta C_p^{\text{fus}}=0$). Здесь $\Phi(T)$ --- идеальная
(кристаллическая) часть, зависящая только от свойств чистого вещества ($\Tm$,
$\dHfus$); $\lng$ --- неидеальность пары.

\paragraph{Модели активности.} Классически $\lng$ задаётся моделью локального
состава. \textbf{NRTL} использует асимметричные параметры $\tau_{12},\tau_{21}$ и
$\alpha$; в пределе бесконечного разбавления
$\ln\gamma_2^\infty=\tau_{12}+\tau_{21}e^{-\alpha\tau_{21}}$. В упрощённом
одно-параметрическом режиме $\lng=\ln\gamma_2^\infty(T)(1-x_2)^2$ с
$\ln\gamma_2^\infty(T)=a+b\,(T_{\text{ref}}/T-1)$ --- аффинно по $1/T$.

\textbf{COSMO-SAC} (используется в нашей новой архитектуре, см.
\S\ref{sec:arch}) выводит $\lng$ из $\sigma$-профилей молекул $p(\sigma)$
(гистограмма экранирующей плотности заряда) через сегментный коэффициент
активности $\Gamma(\sigma)$ (неподвижная точка) и восстановительную свободную
энергию --- физически наиболее обоснованная инженерная модель активности.

\section{Теория: идентифицируемость и условие оптимальности}
\label{sec:theory}
Несущее ядро --- результаты T\ref{thm:ident}, T\ref{thm:cond}, T\ref{thm:temp};
расширения T\ref{thm:semi}, T\ref{thm:geom} и Утв.~\ref{prop:aopt} снабжены явными
предположениями и численной проверкой (\S\ref{sec:diag}).

\begin{assumption}[Активность содержит член $\propto 1/T$]\label{asm:b}
Модель активности имеет свободный параметр при $1/T$ (как $b$ выше или
температурная зависимость NRTL).
\end{assumption}

\begin{theorem}[Структурная неидентифицируемость]\label{thm:ident}
При Предположении~\ref{asm:b} и $\Delta C_p=0$ отображение
$\theta\mapsto\{\mu(T_k;\theta)\}_k$ ранг-дефицитно в направлениях
кристалл/активность при любом наборе температур: чувствительность
кристаллического члена $v_k=\partial\mu(T_k)/\partial\dHfus$ лежит в линейной
оболочке столбцов активности $(\mathbf 1,\,1/T)$. Поэтому $\dHfus$ (и само
расщепление $\Phi/\lng$) \emph{не} идентифицируемо из данных растворимости;
информация Фишера $\mathcal I_{\dHfus}=0$. Дефицит замыкается ровно добавлением
внешнего ограничения на кристаллический фактор (прямая метка $\dHfus$, либо
фиксация наклона активности $b$).
\end{theorem}
\begin{proof}[Схема]
$\Phi(T)=(\dHfus/R)(1/T-1/\Tm)$ даёт лишь форму $\propto 1/T$ плюс константу; обе
уже порождаются членом $b/T$ и константой $a$ активности. Столбец дизайна для
$\dHfus$ --- линейная комбинация существующих, матрица Грама вырождена,
соответствующая информация Фишера равна нулю. Добавление единичной строки по
$\dHfus$ (прямая метка) или удаление столбца $b$ восстанавливает полный ранг.
\end{proof}

\begin{quote}\textbf{Поправка к фольклору.} Мульти-температурные данные сами по
себе \emph{не} идентифицируют расщепление: они лишь добавляют коллинеарные строки
$\propto 1/T$. Это честная версия более ранних оценок вида $(\Delta T)^{-k}$,
которые молча предполагали отсутствие свободного члена $b/T$.\end{quote}

\begin{theorem}[Полупараметрическая граница, класс-зависимая]\label{thm:semi}
Пусть нюисанс-активность лежит в классе $\mathcal H$ с проекцией на касательное
пространство $\Pi_{\mathcal H}$. Эффективная информация о $\dHfus$:
\[ \mathcal I_{\text{eff}}(\dHfus)=\sigma^{-2}\,\big\|v-\Pi_{\mathcal H}v\big\|^2. \]
При богатом классе (свободный $1/T$) $\Pi_{\mathcal H}v=v$ и
$\mathcal I_{\text{eff}}=0$ (согласуется с T\ref{thm:ident}); при ограниченном
классе $\mathcal I_{\text{eff}}>0$ и граница Крамера--Рао конечна.
\end{theorem}

\begin{definition}[Переносимость факторов]
$\rho_c,\rho_a$ --- популяционные риски лучших предсказателей кристаллического и
остаточно-активностного факторов при их покрытиях метками; $\rho_y$ --- риск
прямого предсказателя $\lnx$ при покрытии только пар растворимости.
\end{definition}

\begin{theorem}[Условие оптимальности физики (главный результат)]\label{thm:cond}
На скаффолд-сдвинутом тесте структурированный предсказатель удовлетворяет
\[ \mathcal R_{\text{phys}}^{\text{scaffold}}\le
   \rho_c+\rho_a+\Gamma_{\text{ident}}, \]
где $\Gamma_{\text{ident}}$ --- разрыв идентифицируемости (равен нулю при
замыкании по T\ref{thm:ident}). Физика строго превосходит прямой предсказатель,
$\mathcal R_{\text{phys}}^{\text{scaffold}}<\rho_y$, \emph{тогда и только тогда},
когда (i)~$\Gamma_{\text{ident}}$ замкнут и (ii)~кристаллический фактор наследует
более широкое химическое покрытие, чем пары ($\rho_c$ мал за счёт внешнего
одно-компонентного пула), а остаток активности ниже по сложности, чем совместная
цель.
\end{theorem}
\begin{proof}[Схема]
Разложить структурированный риск через две головы; ограничить каждую её риском
переноса и липшицевостью решателя; неидентифицируемый перекрёстный член собрать в
$\Gamma_{\text{ident}}$. Кристаллический фактор --- одно-компонентная функция,
поэтому супервизируется внешним пулом чистых компонентов много большего покрытия,
чем пары; у прямого предсказателя такого доступа нет.
\end{proof}

\begin{corollary}[Механизм]
Заземление $\Phi$ на внешний пул точек плавления/энтальпий (и GC-приоры)
одновременно замыкает $\Gamma_{\text{ident}}$ и снижает $\rho_c$; симметрично,
заземление активности на внешний пул $\sigma$-профилей снижает $\rho_a$. Оба
рычага реализованы в \S\ref{sec:arch}.
\end{corollary}

\begin{theorem}[Оптимальность температурной экстраполяции]\label{thm:temp}
Для экстраполяции внутри одной пары к $T^\star$ вне диапазона обучения
$[T_{\min},T_{\max}]$ структурированный (аффинный по $1/T$) класс гипотез имеет
ошибку экстраполяции, ограниченную дисперсией оценки наклона, умноженной на
$|1/T^\star-\overline{1/T}|$, тогда как неструктурированный температурный энкодер
не ограничен вне носителя обучения.
\end{theorem}

\begin{theorem}[Инфо-геометрия компенсации, пояснительная]\label{thm:geom}
Минимальное собственное значение метрики Фишера в плоскости кристалл/активность
стремится к нулю при $\Delta T\to0$; значит компенсирующие движения по
поверхности идентификации почти не стоят информации, и оптимизатор их предпочитает.
\end{theorem}

\begin{quote}\textbf{Оговорка о метрике компенсации.} Поскольку
$\delta\Phi+\delta\lng=$ (остаток $\lnx$), корреляция
$\mathrm{corr}(\delta\Phi,\delta\lng)\to-1$ \emph{тривиально} при улучшении
подгонки $\lnx$, независимо от физической корректности. Поэтому робастной метрикой
качества разложения мы берём систематический сдвиг $\overline{\delta\Phi}$, а не
$|\mathrm{corr}|$.\end{quote}

\begin{proposition}[A-оптимальная декомпозиция вспомогательных потерь]\label{prop:aopt}
Среди неотрицательных весов вспомогательных потерь при бюджете trace-оптимальный
выбор даёт положительный вес каждой потере, добавляющей информацию в
ранг-дефицитное направление; при блочной структуре голов (обеспеченной
stop-gradient) это выбирает ровно сигналы на кристалл, на уровень активности и на
форму активности.
\end{proposition}

\section{Диагностика: численная проверка теории}
\label{sec:diag}

\subsection{Аудит идентифицируемости/Фишера (подтверждает T\ref{thm:ident}, T\ref{thm:semi}, T\ref{thm:geom})}
Мы строим матрицу чувствительности SLE-модели и информацию Фишера на
безразмерных параметрах.

\begin{table}[h]
\centering
\caption{Режимы супервизии ($\Delta T=40$\,K): обусловленность $\mathcal I$ и
граница Крамера--Рао на $\dHfus$.}
\begin{tabular}{lrrr}
\toprule
Режим & $\lambda_{\min}(\mathcal I)$ & $\mathrm{cond}(\mathcal I)$ & CRLB sd($\dHfus$), Дж/моль \\
\midrule
только растворимость & $5{,}4\times10^{-2}$ & $2{,}3\times10^{5}$ & $19{,}895$ \\
\;+ метка $\Tm$ & $2{,}32$ & $5{,}4\times10^{3}$ & $3{,}151$ \\
\;+ метки $\Tm$ и $\dHfus$ & $8{,}04$ & $1{,}6\times10^{3}$ & $1{,}689$ \\
\bottomrule
\end{tabular}
\end{table}

\noindent
\textbf{Свип по $\Delta T$} (обусловленность, только растворимость):
$1{,}7\times10^{3}$ при $120$\,K $\to 2{,}3\times10^{5}$ при $40$\,K $\to
1{,}9\times10^{16}$ при $10$\,K --- узкое температурное окно делает расщепление
ранг-дефицитным (T\ref{thm:geom}). \textbf{Класс активности} (вложенное
сравнение): свободный наклон $\Rightarrow$ sd$(\dHfus)=19{,}895$, cond
$2{,}3\times10^5$; известный наклон $\Rightarrow$ sd$(\dHfus)=8{,}251$, cond
$5{,}3\times10^2$ --- идентифицируемость $\dHfus$ \emph{зависит от класса}
активности (T\ref{thm:semi}).

\subsection{Робастный перезамер компенсации}
Строгая метрика компенсации требует измеренных \emph{и} $\Tm$, \emph{и} $\dHfus$
на строке, но в стандартном тест-сплите \textbf{ноль} строк с измеренной
$\dHfus$; ранее публиковавшееся $\mathrm{corr}=-0{,}876$ получено на ручном
наборе из $8$ строк. Относительно физически обоснованного GC-эталона (Joback) на
полном тесте ($n=5826$):
\begin{center}
\begin{tabular}{lr}
\toprule
Величина & Значение \\
\midrule
$\mathrm{corr}(\delta\Phi,\delta\lng)$ & $-0{,}962$ \\
\quad bootstrap CI$_{95}$ & $[-0{,}965;\,-0{,}958]$ \\
$\overline{\delta\Phi}$ (робастная метрика) & $-7{,}1$ \\
доля противоположных знаков & $0{,}93$ \\
\bottomrule
\end{tabular}
\end{center}
Компенсация сильная, устойчивая и отрицательная; кристаллический член
систематически далёк от физического приора, а активность его почти идеально
поглощает --- что мотивирует заземление кристалла (и активности).

\section{Архитектура двойного заземления}
\label{sec:arch}
Логика <<заземлить кристалл внешними одно-компонентными данными>> симметрично
переносится на активность через $\sigma$-профили. Тогда \textbf{оба} фактора
разложения становятся одно-компонентными, внешне-супервизируемыми объектами, и
условие T\ref{thm:cond} выполняется \emph{конструктивно}.

\paragraph{Кристаллическая ветвь.} Голова \textsc{FusionHead} предсказывает
$(\Tm,\dHfus)$ из вектора чистого вещества, заземлённая на GC-приоры (Joback,
вычислимые для любого SMILES) и на внешний пул точек плавления
(\texttt{open\_crystal\_artifact}: $19{,}436$ веществ; после честного исключения
скаффолдов теста/валидации остаётся $14{,}792$, из них $448$ с $\dHfus$). Внешний
пул подаётся отдельным вспомогательным потоком (self-solvent строки, маски только
на $\Tm/\dHfus$), что $\approx\!10\times$ превышает число обучающих веществ с
$\Tm$.

\paragraph{Активностная ветвь (новое).} Сеть предсказывает
\emph{$\sigma$-профиль} молекулы (\textsc{SigmaProfileHead}: softmax по $51$
бину $\times$ softplus-площадь --- одно-компонентный, переносимый объект),
из которого \emph{дифференцируемый} слой COSMO-SAC (\textsc{CosmoSacLayer})
вычисляет $\lng$:
\begin{align}
\Delta W(\sigma_m,\sigma_n) &= \tfrac{\alpha'}{2}(\sigma_m+\sigma_n)^2
   + c_{\text{hb}}\max(0,\sigma_{\text{acc}}-\sigma_{\text{hb}})
     \min(0,\sigma_{\text{don}}+\sigma_{\text{hb}}),\\
\Gamma_S(\sigma_m) &= \Big[\textstyle\sum_n p_S(\sigma_n)\,\Gamma_S(\sigma_n)\,
   e^{-\Delta W(\sigma_m,\sigma_n)/RT}\Big]^{-1}\quad(\text{неподвижная точка}),\\
\ln\gamma_2^{\text{rest}} &= \frac{A_2}{a_{\text{eff}}}\sum_m
   p_2^{\text{pure}}(\sigma_m)\big[\ln\Gamma_S(\sigma_m)-\ln\Gamma_2(\sigma_m)\big],
\end{align}
плюс опциональный комбинаторный член Ставермана--Гуггенхайма. Матрица обмена
$\Delta W$ зависит только от фиксированной сетки и вычисляется один раз; сегментная
неподвижная точка $\Gamma$ решается ограниченной демпфированной итерацией (как
SLE-решатель) и полностью дифференцируема. $\sigma$-профиль супервизируется внешней
базой (VT-2005) отдельным потоком с потерей --- одномерным расстоянием
Вассерштейна (EMD) на форме профиля плюс MSE на площади.

\paragraph{Переключение.} Флаг \texttt{activity\_model="cosmo\_sac"} активирует
$\sigma$-ветвь вместо NRTL; путь NRTL/решателя SLE и контроль \textsc{DirectGNN}
не затронуты. Поправочный механизм корректирует только кристаллическую сторону,
активность проходит без изменений.

\section{Линия ThermoCurve: что сохранено, исправлено и расширено}
\label{sec:thermocurve}
Архитектура двойного заземления (\S\ref{sec:arch}) --- это \textbf{реализованная и
исправленная форма} проектной линии \emph{ThermoCurve}, а не отказ от неё. Здесь
мы явно фиксируем происхождение идей, их математическую запись и то, что именно
сохранено, исправлено и расширено.

\paragraph{Исходный замысел ThermoCurve.} ThermoCurve --- проект архитектуры
второго поколения со тремя принципами: (i)~предсказывать \emph{кривую}
$\widehat{\mathcal C}_{s,v}\colon T\mapsto\widehat{\lnx}(T)$, а не точку;
(ii)~использовать \emph{аналитическую} форму температурной зависимости активности
вместо численного решателя SLE; (iii)~\emph{декомпозировать} функцию потерь так,
чтобы каждая ветвь получала независимый обучающий сигнал. Аналитическая кривая:
\begin{equation}
\widehat{\lnx}(T)=
-\frac{\widehat{\dHfus}}{R}\!\Big(\frac1T-\frac1{\widehat\Tm}\Big)
-\Big[\widehat{\ln\gamma_2^\infty}(T_{\text{ref}})
   +\!\int_{T_{\text{ref}}}^{T}\!\frac{\hat A+\hat B/T'+\hat C\ln T'}{R\,T'^2}\,dT'\Big],
\label{eq:thermocurve}
\end{equation}
где интеграл берётся аналитически (без численного ОДУ-решателя). Декомпозиция
потерь:
\begin{equation}
\mathcal L=\mathcal L_{\text{SLE}}
+\lambda_c\mathcal L_{\text{crystal}}
+\lambda_{\text{IDAC}}\mathcal L_{\text{IDAC}}
+\lambda_{\text{CF}}\mathcal L_{\text{CF}}
+\lambda_{\text{int}}\mathcal L_{\text{integral}}
+\lambda_{\text{mono}}\mathcal L_{\text{mono}},
\label{eq:tc-loss}
\end{equation}
где $\mathcal L_{\text{crystal}}$ --- прямые метки $\Tm,\dHfus$;
$\mathcal L_{\text{IDAC}}$ --- $\ln\gamma_2^\infty$; $\mathcal L_{\text{CF}}$ ---
контрфактический сигнал уровня (для одного вещества в двух растворителях разность
$\lnx$ при $T_{\text{ref}}$ ограничивает разность $\ln\gamma_2^\infty$);
$\mathcal L_{\text{integral}}$ --- чистый сигнал на \emph{форму} кривой через
температурные разности $\lnx(T_k)-\lnx(T_j)$ одной пары; $\mathcal L_{\text{mono}}$
--- мягкая монотонность $\partial\lnx/\partial T\ge0$.

\paragraph{Что сохранено и реализовано.} Все три принципа ThermoCurve --- костяк
нашей архитектуры. Соответствие компонентов:

\begin{center}
\begin{longtable}{p{4.6cm} p{6.3cm} p{2.6cm}}
\toprule
Компонент ThermoCurve & Наша реализация & Статус \\
\midrule
\endhead
Кривая $T\mapsto\lnx$ (не точка) & аналитический активностный путь + SLE-баланс по \eqref{eq:sle} & сохранено \\
Аналитика вместо численного решателя & дифф. активностный слой (замкнутый $\gamma_\infty$ / COSMO-SAC с ограниченной итерацией) + неявное дифференцирование & сохранено \\
$\mathcal L_{\text{crystal}}$ ($\Tm,\dHfus$) & \textsc{FusionHead} + GC-приоры + \textbf{внешний кристалл-пул} (aux-поток) & сохранено и расширено \\
$\mathcal L_{\text{IDAC}}$ & супервизия $\ln\gamma_2^\infty$ + IDAC aux-поток & сохранено \\
$\mathcal L_{\text{integral}}$ (форма) & pair-temperature van't Hoff / $\Delta\lnx$ + exact-joint supervision & сохранено (частично) \\
$\mathcal L_{\text{CF}}$ (уровень) & ранжирование по T внутри пары + декорр-диагностика & частично \\
$\mathcal L_{\text{mono}}$ & монотонность по $T$ (фаза 3) & сохранено \\
Форма активности $(A,B,C)$/Лежандр & \textbf{заменена} на $\sigma$-профиль$\to$COSMO-SAC & расширено \\
Башня CRLB-теорем & ранговая T\ref{thm:ident} + полупараметрика T\ref{thm:semi} + условие T\ref{thm:cond} & исправлено \\
$\mathrm{corr}=-0{,}876$ на $8$ строках & робастный полнотестовый GC-замер ($-0{,}962$, с оговоркой) & исправлено \\
\bottomrule
\end{longtable}
\end{center}

\paragraph{Что исправлено.} В исходном предложении ThermoCurve теоретическая
часть была внутренне противоречива: конечная граница Крамера--Рао (его Теорема
1.1) противоречила собственному следствию о нулевой эффективной информации;
содержала ошибку в линейной оболочке ($v=u-1/\Tm\in\operatorname{span}(1,u)$);
а оценка $(\Delta T)^{-4}$ молча отбрасывала свободный член $B/T$ активности.
Мы заменили эту башню корректным ранговым результатом T\ref{thm:ident},
класс-зависимой полупараметрической границей T\ref{thm:semi} и условием
оптимальности T\ref{thm:cond} (\S\ref{sec:theory}). Эмпирику компенсации на $8$
строках ($\mathrm{corr}=-0{,}876$) мы заменили устойчивым полнотестовым замером
относительно GC-эталона ($-0{,}962$, $n=5826$), с честной оговоркой о
тривиальности самой корреляции (\S\ref{sec:diag}).

\paragraph{Что расширено (двойное заземление).} ThermoCurve заземлял только
кристаллическую ветвь; его форма активности $(A,B,C)$ \emph{не имела прямых меток}
и обучалась лишь косвенно через интегральный/разностный сигнал. Мы сделали
активностную ветвь \emph{тоже} внешне-заземлённой: сеть предсказывает
$\sigma$-профиль (одно-компонентный, переносимый объект), супервизируемый внешней
базой, который подаётся в дифференцируемый COSMO-SAC. В результате \textbf{оба}
фактора разложения становятся одно-компонентными внешне-супервизируемыми
объектами --- это строго сильнее одиночного (кристаллического) заземления
ThermoCurve и есть конструктивная реализация условия T\ref{thm:cond}: COSMO-SAC
физичнее трёх-параметрической формы и сам несёт правильную температурную
зависимость через $\Delta W/RT$.

\begin{quote}\textbf{Вывод по линии ThermoCurve.} Мы её не отбросили: её ядро
(кривая, аналитика, раздельная супервизия) сохранено и составляет костяк
архитектуры; слабые места (противоречивая CRLB-теория, $8$-строчная эмпирика,
бедная параметризация активности) исправлены; а само заземление расширено с
одного фактора до двух.\end{quote}

\section{Оценка по утилите (а не только MAE)}
\label{sec:utility}
Абсолютный MAE по $\lnx$ --- неверная метрика для выбора растворителя и
чувствителен к шуму и систематическим сдвигам. Мы добавили оценку, отражающую
реальную задачу.

\paragraph{Ранжирование растворителей.} Для каждой группы $(\text{вещество},T)$
с $\geq3$ растворителями ранжируем растворители по предсказанной и измеренной
$\lnx$ (инвариантно к сдвигам и шуму меток). Прокси-TGNN на scaffold-тесте
($826$ групп, в среднем $6{,}3$ растворителя):
\begin{center}
\begin{tabular}{lrrr}
\toprule
$\rho_{\text{Spearman}}$ & Kendall $\tau$ & top-1 (лучший растворитель) & NDCG@3 \\
\midrule
$0{,}51$ & $0{,}44$ & $0{,}48$ & $0{,}73$ \\
\bottomrule
\end{tabular}
\end{center}
То есть при <<плохом>> MAE $1{,}74$ модель выбирает лучший растворитель в $48\%$
случаев (против $\approx16\%$ случайно) --- практическая полезность много выше,
чем подсказывает MAE.

\paragraph{Калибровка неопределённости.} Многоуровневый PICP + регрессионная ECE
+ <<резкость>> с учётом шумового потолка (MC-dropout). Прокси-TGNN сильно
\textbf{переуверен}: при номинальном покрытии $90\%$ реальное --- лишь $0{,}13$
(ECE $=0{,}69$; средний std $0{,}28$, что $1{,}7\times$ шумового потолка). Штатная
проверка только на $90\%$ недооценила бы это. \textbf{Решение реализовано}:
split-конформное предсказание даёт \emph{гарантированное} маргинальное покрытие
независимо от того, насколько неверна собственная $\sigma$ модели --- на том же
прокси-TGNN номинал $90\%$ даёт эмпирические $0{,}918$ (ширина $\pm1{,}5\
\log_{10}$, что честно отражает трудную scaffold-неопределённость; адаптивный
нормированный вариант сужает интервалы там, где модель уверена). Вывод:
исходная неопределённость переуверена, но конформ делает её \emph{валидной для
применения}.

\section{Сравнение с внешними моделями (FastSolv, SolProp)}
\label{sec:external}
Прямое сравнение требует осторожности: внешние модели предсказывают $\log_{10}S$
в \emph{молярности} (моль/л), а мы --- $\lnx$ в \emph{мольной доле}. Пересчёт:
$\ln\to\log_{10}$ делением на $\ln 10=2{,}303$; переход мольная доля
$\to$ молярность --- аддитивный сдвиг $\approx{+}1$ на растворитель (влияет на
уровень, почти не на разброс). Поэтому RMSE сопоставимы с точностью
$\pm0{,}1$--$0{,}2\ \log_{10}$. Наш TGNN: RMSE $2{,}34$ нат-лог
$\to 1{,}02\ \log_{10}$; MAE $1{,}74\to 0{,}76\ \log_{10}$.

\begin{table}[h]
\centering
\caption{Строгая экстраполяция на \emph{новые вещества}, RMSE в $\log_{10}$
(органика, T-зависимость). Внешние числа --- из FastSolv (Nature Comm.\ 2025)
и SolProp/Vermeire (JACS 2022).}
\begin{tabular}{lll r}
\toprule
Модель & Тип & Тест & RMSE $\log_{10}$ \\
\midrule
\textbf{FastSolv} (SOTA, 2025) & прямой/чёрный ящик & SolProp (var-$T$) & $0{,}83$ \\
\textbf{FastSolv} & прямой/чёрный ящик & Leeds (room-$T$) & $0{,}95$ \\
\textbf{Наша} (TGNN/DirectGNN) & физ-информ. & BigSolDB scaffold & $\approx 1{,}0$ \\
SolProp/Vermeire (2022) & физ-информ. (термо-цикл) & SolProp & $1{,}43$ \\
SolProp/Vermeire & физ-информ. & Leeds & $2{,}16$ \\
\midrule
\textit{(справка)} FastSolv & интерполяция & BigSolDB train/val & $0{,}22$ \\
\bottomrule
\end{tabular}
\end{table}

\paragraph{Вердикт.} На сопоставимом режиме (новые вещества, органика,
T-зависимость) наш $\approx1{,}0\ \log_{10}$: \emph{немного хуже} текущего SOTA
FastSolv ($0{,}83$--$0{,}95$) и \emph{заметно лучше} прежнего SOTA
SolProp/Vermeire ($1{,}43$--$2{,}16$). Мы в той же лиге, что современные модели,
чуть позади лидера. (Оговорка: тест-сеты не совпадают --- наш scaffold-сплит
BigSolDB против их внешних Leeds/SolProp; пересчёт единиц приблизителен.)

\begin{quote}\textbf{Ключевое наблюдение.} SOTA-модель FastSolv --- это
\emph{прямой чёрный ящик} (дескрипторы fastprop $+$ ансамбли $+$ очищенный
BigSolDB~2.0), \emph{без} термодинамического узкого места; а физ-информированный
SolProp --- \emph{позади} неё. Это внешне подтверждает наш внутренний вывод:
физический bottleneck не выигрывает в сырой точности, а рычаг --- представление,
данные и ансамбли, а не топология энкодера. Отсюда честная ниша нашей работы ---
не <<побить MAE>>, а \emph{конкурентная с SOTA, но интерпретируемая и заземлённая
на оба фактора} модель, с акцентом на температурную экстраполяцию (FastSolv
плато́ит у ${\sim}350$\,K, где принципиальная физика $T$-зависимости --- наше
преимущество).\end{quote}

\section{Текущие результаты (сводка)}
\begin{table}[h]
\centering
\caption{Опорные числа (scaffold-сплит, если не указано иное).}
\begin{tabular}{lrr l}
\toprule
Модель / метрика & MAE & $R^2$ & Примечание \\
\midrule
\textsc{DirectGNN} & $1{,}652$ & $0{,}478$ & лучший по MAE \\
RF (гибрид) & $1{,}722$ & $0{,}449$ & сильный дескрипторный бейзлайн \\
\textsc{TGNN} (MPNN) & $1{,}741$ & $0{,}438$ & физика, текущий <<налог>> $+0{,}09$ \\
RF (дескрипторы) & $1{,}20$ & --- & узкое место = представление \\
pair van't Hoff (T-экстрап.) & $0{,}368$ & $0{,}887$ & оракул структуры по $1/T$ \\
\midrule
Шумовой потолок $\lnx$ & \multicolumn{3}{l}{медиана $0{,}031$ / среднее $0{,}161$ (\S\ref{sec:noise})} \\
Ранжирование (TGNN) & \multicolumn{3}{l}{Spearman $0{,}51$, top-1 $0{,}48$ (\S\ref{sec:utility})} \\
Калибровка (TGNN) & \multicolumn{3}{l}{PICP$_{90}=0{,}13$, ECE $0{,}69$ (\S\ref{sec:utility})} \\
\bottomrule
\end{tabular}
\end{table}

\section{Что реализовано и проверено}
Все локально-выполнимые компоненты реализованы и покрыты тестами (полный набор:
$284$ теста зелёные):
\begin{itemize}
\item \textbf{Phase A}: скрипт оценки шумового потолка (сырой BigSolDB).
\item \textbf{Phase B1}: дифференцируемый \textsc{CosmoSacLayer} (сегментная
  неподвижная точка, предвычисленная $\Delta W$, restoring + комбинаторный член),
  подключён через \texttt{activity\_model="cosmo\_sac"} к решателю SLE.
\item \textbf{Phase B2}: \textsc{SigmaProfileHead} и переключение голов в модели.
\item \textbf{Phase B3}: внешняя $\sigma$-супервизия --- builder потока, ingestion
  VT-2005, потеря EMD, хук в тренере, CLI; полный прогон \texttt{cosmo\_sac}+$\sigma$
  работает end-to-end.
\item \textbf{Phase C}: ранжирование (Spearman/Kendall/top-1/NDCG) и калибровка
  неопределённости (PICP/ECE, noise-floor-aware).
\end{itemize}

\section{Ограничения и дальнейшая работа}
\begin{itemize}
\item Сырой scaffold-MAE ограничен энкодером ($\sim\!93\%$ разрыва до RF ---
  представление, не физический слой); мы \emph{не} заявляем SOTA по MAE.
\item Per-molecule $\dHfus$ слабо покрыт метками ($\sim\!2\%$); опора на GC-приоры
  и энтропийное (Walden) ограничение.
\item GC-эталон в диагностике компенсации --- прокси, не истина; робастен знак и
  относительный $\overline{\delta\Phi}$, не абсолют.
\item $\mathrm{corr}(\delta\Phi,\delta\lng)\to-1$ тривиально улучшается с
  подгонкой $\lnx$; отчётная метрика --- $\overline{\delta\Phi}$.
\item \textbf{Нужны данные/вычисления}: реальная база $\sigma$-профилей (VT-2005;
  \texttt{opencosmorspy} её не поставляет) для настоящего покрытия супервизии;
  GPU для полноценного прогона \texttt{cosmo\_sac} vs \texttt{gamma\_inf} и
  скаффолд/температурных протоколов.
\item Неопределённость требует рекалибровки (temperature scaling / ансамбли).
\end{itemize}

\section{Заключение}
Физика помогает обобщению ровно тогда, когда её факторизация идентифицируема, а
факторы переносимее совместной цели. Практическое следствие --- заземлять оба
фактора на внешние одно-компонентные данные (точки плавления для кристалла,
$\sigma$-профили для активности). Шумовой потолок показывает, что точность
выигрываема; диагностика подтверждает неидентифицируемость и её замыкание внешними
метками; оценка по утилите показывает, что модель уже полезна для ранжирования,
хотя её неопределённость требует калибровки. Архитектура двойного заземления
построена и протестирована; остаются данные ($\sigma$-база) и вычисления (GPU).

\end{document}
